\clearpage
\VerdierPhysicalPageAnchor{0202}%
\begingroup\small\setlength{\parskip}{0pt}\setlength{\emergencystretch}{1.5em}
\setlength{\abovedisplayskip}{4pt plus 2pt minus 2pt}%
\setlength{\belowdisplayskip}{4pt plus 2pt minus 2pt}%

\VerdierSemanticAnchor{ch3d-2-3-3-proof-q-r-step}
\[
\alpha\colon L^{-r+1}\longrightarrow Z^{-r+1},
\qquad\text{ここで }L^{-r+1}\text{ は }M\text{ の対象である}.
\]
次をおく：
$X_{q,r}=X_{q,r-1}(\alpha,-r+1)$、
$v_{q,r}=p(\alpha,-r+1)$ (2.1.1)。$X_{q,r}$ と $v_{q,r}$ が
性質 $I')_r$、$III')_r$、$IV')_r$ をもつことは明らかである。
さらに、複体 $X_{q,r}$ の定義 (2.1.1) により、次の完全列がある：
\[
0\longrightarrow (X_{q,r})^{-r}\longrightarrow
L^{-r}\oplus L^{-r+1}\longrightarrow Z^{-r+1}\longrightarrow0.
\]
したがって、性質 $(E_4)$ から $(X_{q,r})^{-r}$ は $M'$ に属し、
(2.2.6) から
\[
M\text{-}\dim((X_{q,r})^{-r})
=M\text{-}\dim(Z^{-r+1})-1
=M\text{-}\dim(Y^q)-r-q,
\]
すなわち性質 $II')_r$ が従う。

\VerdierSemanticAnchor{ch3d-2-3-3-proof-q-completion}
よって構成を続けることができる。
$X_{q,-q+M\text{-}\dim(Y^q)}=X_q$ とおき、
$u_q\colon X_q\longrightarrow X_{q-1}$ を $v_{q,l}$ の合成で得られる射とする。
$X_q$ と $u_q$ が性質 $1)_q$、$2)_q$、$3)_q$、$4)_q$、$5)_q$ を
もつことは明らかである。これで命題の第一の主張の証明が完了した。
第二の主張を証明するには、付加的な仮定のもとで、適当な $n_0$ に対して
上の構成で $X_{n_0}=Y$ ととれることを注意すれば十分である。

\VerdierSemanticAnchor{ch3d-2-3-4}
\medskip\noindent\textbf{注意 2.3.4.}\quad
複体 $X_{n_0}$ の構成（命題 (2.3.3) の証明を参照）において、実際には
次の結果を証明した。$M\subset M'$ を $\mathcal A$ の対象の二つの集合とし、
$0\in M$、$M'$ は性質 $(E_1)$ と $(E_4)$ をもち、$M'$ の各対象は
$M$ のある対象の商対象であるとする。
$Y\in\operatorname{Ob}\operatorname{comp}^{-}(\mathcal A)$ を、すべての成分が
$M'$ に属する複体とする。このとき、すべての成分が $M$ に属する複体
$X\in\operatorname{Ob}\operatorname{comp}^{-}(\mathcal A)$ と擬同型
$u\colon X\longrightarrow Y$ が存在する。

\VerdierSemanticAnchor{ch3d-2-3-5}
\medskip\noindent\textbf{2.3.5.}\quad
この節の結果を $\mathcal A$ の反対圏に適用して得られる主張の定式化は
読者に委ねる。$\mathcal A$ の対象の集合 $M$ が性質 $(E_i)^o$、
$1\leq i\leq6$ をもつとは、$\mathcal A^o$ の対象の集合 $M$ が
(2.2.1) の性質 $(E_i)$ をもつことをいう。

\VerdierSemanticAnchor{ch3d-2-4}
\subsection*{2.4. 応用。}

\VerdierSemanticAnchor{ch3d-2-4-1}
\noindent\textbf{命題 2.4.1.}\quad
\textit{a) $\mathcal A$ をアーベル圏とし、$M$ を性質 $(E_1)$ (2.2.1) をもつ
$\mathcal A$ の対象の集合とする。圏 $K^{-}(M)$ (1.1.4) が左
$\operatorname{Ac}(\mathcal A)$-局所化部分圏（第 II 章，2.3.5）であるための
必要十分条件は、集合 $M$ が性質 $(E_2)$ と $(E_3)$ をもつことである。}

\noindent\textit{b) $M$ を性質 $(E_1)$、$(E_2)$、$(E_3)$ をもつ対象の集合とし、
$M^{\wedge}$ を (2.2.1) で定義された $M$ に付随する集合とする。自然関手}
\[
K^{-}(M)/\operatorname{Ac}^{-}(M)
\longrightarrow
K^{-}(M^{\wedge})/\operatorname{Ac}^{-}(M^{\wedge})
\]
\endgroup
